\section{Split-prime and graph certificate}
\label{app:split-graph}

This appendix records the finite objects used in the proofs of
\cref{lem:embedded-subgroups,lem:split-noncollision,lem:graph-bridge}.  All
indices are one-based and refer to the fixed
line labelling.  Thus the lists below, together with the cycle arrays in
\cref{eq:Hplus-cycles,eq:Hminus-cycles}, are independent of subgroup-table
positions.

\subsection{The two supports}

The first orbit in $\cS_+$ is
\begin{equation*}
\begin{aligned}
\cO_{+,1}=\{&\{1,2\},\{1,18\},\{2,18\},\{3,6\},\{3,19\},\{4,5\},\\
&\{4,7\},\{5,7\},\{6,19\},\{8,11\},\{8,23\},\{9,13\},\\
&\{9,26\},\{10,14\},\{10,25\},\{11,23\},\{12,15\},\{12,22\},\\
&\{13,26\},\{14,25\},\{15,22\},\{16,17\},\{16,27\},\{17,27\},\\
&\{20,21\},\{20,24\},\{21,24\}\}.
\end{aligned}
\end{equation*}
The second orbit is
\begin{equation*}
\begin{aligned}
\cO_{+,2}=\{&\{1,9\},\{1,25\},\{2,14\},\{2,26\},\{3,12\},\{3,23\},\\
&\{4,16\},\{4,24\},\{5,21\},\{5,27\},\{6,11\},\{6,22\},\\
&\{7,17\},\{7,20\},\{8,15\},\{8,19\},\{9,25\},\{10,13\},\\
&\{10,18\},\{11,22\},\{12,23\},\{13,18\},\{14,26\},\{15,19\},\\
&\{16,24\},\{17,20\},\{21,27\}\}.
\end{aligned}
\end{equation*}
Thus $\cS_+=\cO_{+,1}\sqcup\cO_{+,2}$ has $54$ pairs.  Its family-list,
stabilizer-inventory, value-list, and sorted-value digests are respectively
\begin{align*}
&\hexsha{0b2fa74280691fd83d364d150ff187431b548b8e9e8c951e6dbae034c189c2ec},\\
&\hexsha{500f61f9101ca424c9ff142526c26e41fc2b06c1478b7f2feaf626f07c28aafe},\\
&\hexsha{950b67aa2693ddf9608a22e90348c0a7dfdb71942611fa8ff4fbf392c649afa7},\\
&\hexsha{289d4fcdd86554490a4208ea4fb8ead6abcd1e1d3b8b162c6e2d3805dd0fcee5}.
\end{align*}

The negative support is the following single $81$-element orbit:
\begingroup\small
\begin{equation*}
\begin{aligned}
\cS_-=\{&\{1,2\},\{1,5\},\{1,6\},\{1,13\},\{1,19\},\{1,20\},\\
&\{2,3\},\{2,4\},\{2,11\},\{2,18\},\{2,20\},\{3,5\},\\
&\{3,6\},\{3,7\},\{3,10\},\{3,11\},\{4,5\},\{4,6\},\\
&\{4,12\},\{4,17\},\{4,18\},\{5,7\},\{5,12\},\{5,13\},\\
&\{6,10\},\{6,17\},\{6,19\},\{7,9\},\{7,15\},\{7,24\},\\
&\{7,26\},\{8,12\},\{8,17\},\{8,18\},\{8,23\},\{8,26\},\\
&\{8,27\},\{9,10\},\{9,11\},\{9,23\},\{9,26\},\{9,27\},\\
&\{10,14\},\{10,24\},\{10,27\},\{11,21\},\{11,23\},\{11,24\},\\
&\{12,15\},\{12,25\},\{12,26\},\{13,15\},\{13,16\},\{13,22\},\\
&\{13,26\},\{14,17\},\{14,19\},\{14,22\},\{14,24\},\{14,25\},\\
&\{15,22\},\{15,24\},\{15,25\},\{16,19\},\{16,20\},\{16,23\},\\
&\{16,26\},\{16,27\},\{17,25\},\{17,27\},\{18,21\},\{18,23\},\\
&\{18,25\},\{19,22\},\{19,27\},\{20,21\},\{20,22\},\{20,23\},\\
&\{21,22\},\{21,24\},\{21,25\}\}.
\end{aligned}
\end{equation*}
\endgroup
Its corresponding four digests are
\begin{align*}
&\hexsha{4e329548ae15632deabfa8c1bea11b91d2fc7c60ddfc86ee4110d86a54d9cb66},\\
&\hexsha{11b388fe39f5e884f064c38fda924cd14e9d063736cedee16d2b9cc01b3c1996},\\
&\hexsha{aaffea3db8c0c913517562f702d07ab9f1b86c78ee90821d7aa09984f6d239c7},\\
&\hexsha{b45cbd50b4a2c278416b6dd4a83ac7b160ad5113ff1b7d2a96e6c8961c0ee36b}.
\end{align*}

\subsection{Finite-field noncollision}

At the certified split prime $p=692717$, the eliminant has $27$ distinct
roots in $\F_{692717}$.  The ordered root vector has digest
\[
 \hexsha{96c072b9f1cc6a3e07d28433ed4b864b50f7d6841ed84a6536fd43ee050017dd}.
\]
For each sign, evaluating the scaled support sum on all $320$ cosets gives
$320$ distinct residues.  The preceding support digests and the
root-vector digest bind this assertion to the exact computation; no
decimal approximation enters.  Distinctness modulo $692717$ proves
distinctness in characteristic zero, as used in
\cref{lem:split-noncollision}.

\subsection{The graph authority chain}

The certified line-coordinate array has digest
\[
 \hexsha{bedb9a147cc2d196ab00b91bc147f6d544b8697fc9458ed269c0e52671a62dfd}.
\]
The reconstructed edge list and the independently generated standard
line-intersection edge list have digests
\begin{align*}
 &\hexsha{c57dc120a33de30dc506e5ab045152b7d55dc2876685621ff21dad99586ffa8f},\\
 &\hexsha{c10e51e18748e3474a6338f612a38e28c74eee61d24922d6626aa1de04bbaa1b}.
\end{align*}
There are $135$ edges, every vertex has degree $10$, and the two graphs
are isomorphic.  This is the complement of the conventional skew-line
Schlaefli graph, which is $16$-regular; the two conventions have the same
automorphism group.  The chosen isomorphism has digest
\[
 \hexsha{a986674a4a0e64dc54244b81a491b86b4fe0af252fd5ad86512832641a64b0ea}.
\]
The transported standard generators have digest
\[
 \hexsha{e91d1484d18cca37378c5f8adbb5a74e897b0befc07a17f4cafbbbad8a99bcc1},
\]
while the enumerated $W(E_6)$ element list has digest
\[
 \hexsha{2ad218397767b236bd90c78ab3a8cd8f8b791d838efd8f9be57fab9c14d1839d}.
\]
The upper-bound stabilizer cells have sizes $27,16,10,6,2$, whose product
is $51840$.  The enumerated subgroup has that same order, so the graph
automorphism group and the labelled Galois image coincide exactly.
